\chapter{CONCLUSÃO}
\label{ch:conclusao}

Este trabalho investigou como o tipo e a ordem das perspectivas atribuídas a agentes baseados em grandes modelos de linguagem se associam à transformação do conteúdo e do enquadramento emocional de notícias em cadeias de reescrita. Para observar esse processo, foi desenvolvido o \textit{Structured Topic Drift Index} (STDI), examinado por comparações quantitativas, casos controlados e inspeções qualitativas. O estudo reuniu três geradores e dois corpora de 50 entradas cada, com oito trajetórias de quatro posições.

Os resultados apresentados no Capítulo \ref{ch:resultados} e discutidos no Capítulo \ref{ch:discussao} permitem uma resposta mais definida sobre os contrastes de perspectiva do que sobre a ordem das reescritas. A introdução da perspectiva conspiratória após duas passagens neutras esteve associada a maior STDI final que a continuação inteiramente neutra, nas 50 entradas de cada uma das seis rodadas. Como as duas continuações partem da mesma mensagem no prefixo compartilhado, o contraste evidencia a diferença entre os resultados das instruções de preservação e de interpretação conspiratória nas condições estudadas.

As comparações de sequências com a mesma composição de perspectivas, por sua vez, não apresentaram uma direção consistente entre geradores e corpora. Isso não demonstra que a ordem seja irrelevante. A inversão também altera a perspectiva responsável pela última reescrita, impedindo separar sua contribuição daquela das posições anteriores. Os resultados descrevem, portanto, diferenças entre percursos completos, sem estabelecer um efeito isolado da ordem.

O acompanhamento das versões intermediárias mostrou ainda que preservar a mensagem recebida pode manter uma transformação anterior. Quando a retransmissão neutra sucedeu a perspectiva conspiratória, pequenos incrementos coexistiram com afastamento elevado do original. O agente neutro conservou a interpretação disponível, conforme sua instrução, sem restaurar automaticamente a notícia de origem. Esse comportamento torna necessária a leitura conjunta da mudança local e da distância ao original.

Essa transformação também pôde ocorrer com a preservação do acontecimento central e de seus participantes. Subtópicos e relações apresentaram as maiores médias de deriva entre os componentes de conteúdo, e a leitura dos casos revelou alterações de motivos, justificativas e responsabilização. Nas trajetórias com a perspectiva conspiratória, o VAD por LLM indicou, em média, menor valência e maior ativação em relação ao original. Esses sinais descrevem o tom expresso e complementam a análise qualitativa do enquadramento, sem esgotá-la.

Para tornar esse percurso examinável, o trabalho oferece um procedimento documentado de geração e comparação sobre as mesmas entradas, com preservação de textos, estruturas, configurações e escores. Os prefixos compartilhados permitem comparar continuações de uma mesma mensagem, enquanto o pareamento por notícia mantém explícitas as dependências do desenho. A decomposição do STDI em tema, subtópicos, entidades e relações, acompanhada de contradição interna e VAD, permite localizar dimensões que um escore final isolado não explica. Os registros preservados permitem recompor as análises realizadas.

As comparações quantitativas e a inspeção dos textos também delimitaram o alcance do instrumento. O desacordo entre Cluster e LLM persistiu com VAD comum, e parte das divergências foi relacionada a inconsistências da extração. Houve ainda oscilação de escore entre versões equivalentes após normalização de espaços. Esses diagnósticos mostram que o valor observado depende tanto da transformação textual quanto da representação e do julgamento utilizados para medi-la. Uma redução do STDI exige inspeção antes de ser interpretada como recuperação informacional.

As contribuições têm caráter instrumental e descritivo. A ausência de referência humana numérica independente limita a avaliação quantitativa do STDI, e as diferenças entre corpora e configurações de mensuração impedem atribuir seus contrastes apenas ao idioma. A análise conjunta de trajetórias, componentes e textos permite examinar onde a mensagem muda e como interpretações acrescentadas permanecem nas transmissões seguintes.

\section{TRABALHOS FUTUROS}
\label{sec:conclusao-futuros}

Uma continuidade metodológica seria comparar trajetórias com a mesma composição e a mesma perspectiva final, variando as posições anteriores, para examinar a ordem com menor sobreposição entre fatores. Repetições independentes da geração e maior controle dos parâmetros de amostragem ajudariam a distinguir padrões recorrentes de particularidades das realizações preservadas.

No instrumento, caberia examinar a estabilidade da extração e do julgamento diante de variações textuais equivalentes, distinguindo resultados reutilizados por cache de novas avaliações. Comparações com outros avaliadores, variações de prompt e análises de sensibilidade dos pesos também poderiam esclarecer quanto as conclusões dependem das escolhas operacionais.

Uma camada complementar de verificação factual externa permitiria investigar quando a transformação textual é acompanhada de alteração na correção das afirmações, com fontes e critérios próprios. A ampliação para novos corpora, outros idiomas e modelos além dos três estudados permitiria examinar a persistência dos padrões. Uma investigação de diferenças linguísticas exigiria controlar o conteúdo entre idiomas, pois os corpora atuais não são traduções das mesmas notícias.

Por fim, redes com múltiplas fontes, decisões de exposição e compartilhamento ou intervenções de correção ampliariam o objeto para além da retransmissão entre antecessores únicos. Essa extensão exigiria procedimentos próprios para distinguir a transformação da mensagem dos comportamentos atribuídos aos agentes.
